% Source-aligned English for printed p.193.
\clearpage
\SourcePage{193}

\LayerTarget{tohoku:remark:4.3.2:1}
still true, however (for other reasons), if $\mathrm C(A)$ is replaced by the
``Cartan resolution'' $A\otimes_ZF$, where $F$ is a ``fundamental
sheaf'' on $X$, assumed paracompact.

\LayerTarget{tohoku:remark:4.3.2:2}
2. Let $x\in X$. We saw in 4.2 that
\LayerTarget{tohoku:formula:stalk-ext-filtered-colimit}
$\mathbf{Ext}_{\mathbf O}^p(\mathbf M,A)(x)=
\mathop{\underset{\to}{\mathrm{lim.}}}\mathrm{Ext}_{\mathbf O}^p(U;\mathbf M,A)$, where the inductive limit is
taken over the directed set of open neighborhoods $U$ of $x$.
\LayerTarget{tohoku:construction:stalk-spectral-sequence-limits}
For each $U$ one may form spectral sequences $II$ and $III$. Passing readily
to the inductive limit shows that
$(\mathbf{Ext}_{\mathbf O}^p(\mathbf M,A)(x))_p$ \textit{is the abutment of two
spectral sequences} whose initial terms are, respectively,

\LayerTarget{tohoku:formula:stalk-spectral-ii-initial-term}
\[
II_2^{p,q}=\mathop{\underset{\to}{\mathrm{lim.}}}
\mathrm{Ext}_O^p(M,H^q(U,A))
\]

\LayerTarget{tohoku:formula:stalk-spectral-iii-initial-term}
\[
III_2^{p,q}=
\bigoplus_{q'+q''=q}\mathop{\underset{\to}{\mathrm{lim.}}}
\mathrm{Ext}_K^p(\mathrm{Tor}_{q'}^O(K,M),H^{q''}(U,A)).
\]
(In the second spectral sequence, $A$ is assumed to be annihilated by a
two-sided ideal $I$, and $K=O/I$.)
\LayerTarget{tohoku:criterion:stalk-ext-limit-transfer}
It follows that whenever, in one of these spectral sequences, the $\mathop{\underset{\to}{\mathrm{lim.}}}$ sign can be transferred to $H(U,A)$ in each component of the initial term, the
formula
$\mathbf{Ext}_{\mathbf O}^p(\mathbf M,A)(x)=
\mathrm{Ext}_O^p(M,A(x))$ holds. Its validity was previously assured only
when $M$ admits a projective resolution by finite free modules $O^{n_i}$.
Here are two typical cases in which this occurs:
\LayerTarget{tohoku:example:constant-sheaf-hlc}
a) $A$ is the \textit{constant} sheaf defined by an $O$-module $N$, and $X$ is HLC
(the first spectral sequence is used).
\LayerTarget{tohoku:example:noetherian-quotient}
b) $A$ is annihilated by an ideal $I$ such that $O/I=K$ is noetherian, and the
$\mathrm{Tor}_q^O(K,M)$ are of finite type over $K$. ---Suppose that $A$ is
the constant sheaf defined by an $O$-module $N$ annihilated by $I$, and, for
simplicity, suppose that $K=O/I$ is a field. Then
\LayerTarget{tohoku:formula:constant-sheaf-ext-presheaf}
$\mathbf{Ext}_{\mathbf O}^n(\mathbf M,\mathbf N)$ is the sheaf associated to
the presheaf defined by
$\displaystyle \bigoplus_{p+q=n}
\mathrm{Hom}_K(\mathrm{Tor}_p^O(K,M),H^q(U,N)).$
\LayerTarget{tohoku:example:theorem-4.2.2-stalk-failure}
It is easy to construct examples (with $M=N=K$, $p=1$, $O$ the group algebra
$K(G)$ of a group $G$, and $I$ the augmentation ideal) in which this sheaf is
distinct from the constant sheaf associated to
$\mathrm{Ext}_O^p(M,N)$ ($=H^p(G,K)$). In this case the conclusion of Theorem
4.2.2 therefore fails.

\LayerTarget{tohoku:section:4.4}
\textbf{4.4. Case of sheaves with a group of operators.}
\LayerTarget{tohoku:setup:group-operator-sheaves}
Let $G$ be a group and $k$ a commutative ring with identity, and put
$O=k(G)$, the group algebra of $G$ with coefficients in $k$. Then $O$ is an
augmented algebra; denote its augmentation ideal by $I$, so that $k$ is
identified with the $O$-module $O/I$. An $\mathbf O$-module $A$ is simply a sheaf of
$k$-modules admitting $G$ as a group of operators. To say that $A$ is
annihilated by $I$ means that $G$ acts trivially on $A$. Because of their
importance for the next chapter, we shall quickly review the essential
notation and results of the preceding sections in the present case.
\LayerTarget{tohoku:setup:g-module-sheaf-notation}
Suppressing $k$ once and for all from the notation (in practice, $k$ will be
the ring $\mathbf Z$ of integers, or a field), we shall say $G$-module or
$G$-sheaf in place of $\mathbf O$-module. If $A,B$ are two $G$-sheaves, write
$\mathrm{Hom}_G(B,A)$, $\mathbf{Hom}_G(B,A)$,
$\mathrm{Ext}_G^p(X;B,A)$, and $\mathbf{Ext}_G^p(B,A)$, putting $G$ in the
subscript in place of $\mathbf O$. When $B=k$, the only case of importance below,
the preceding objects will also be denoted $\Gamma^G(A)$, $A^G$,
$H^p(X;G,A)$, and

\TohokuContinuation{194}
